The novel part of the proof of \Cref{thm:app-main} is the following comparison theorem which allows us to relate vanishing lines for different rings.

\begin{thm}[Comparison Theorem]
  \label{thm:AB-vl}\
  
  Suppose we have a ring map $i : A \to B$, then
  \begin{enumerate}
  \item $g_A(k) \leq g_B(k)$,
  \item $f_A(k) \leq f_B(k)$,
  \item if $i$ becomes an equivalence after applying $\tau_{<m}$, then
    $$ f_B(k) \leq f_A(k) + \left\lfloor \frac{k + f_A(k) - 1}{m} \right\rfloor \leq \frac{1}{m}k + \left( 1 + \frac{1}{m} \right) f_A(k) - \frac{1}{m}. $$
  \end{enumerate}
\end{thm}

\begin{rmk}
  In \cite{BBBCX2}, Conjecture 9.4.2 asks whether there is a finite-page vanishing line of slope $\frac{1}{13}$ in the $\mathrm{tmf}$-Adams spectral sequence for a particular spectrum. We can provide the following evidence in favor of this conjecture:
  The map
  $$ \mathrm{tmf} \to \mathrm{tmf}_1(3) = \BP\langle 2 \rangle $$
  allows us to apply \Cref{thm:AB-vl}(2), \Cref{thm:app-main}(1) and the Nilpotence theorem in order to conclude that
  $$ f_{\mathrm{tmf}}(k) \leq f_{\BP\langle 2 \rangle}(k) \leq \frac{1}{14}k + o(k) $$ 
  Note that the bound on $f_{\mathrm{tmf}}$ is not guaranteed to appear at any finite page. \todo{Perhaps note the consequence for the $v_2 ^2$-elements in the telescope?}
\end{rmk}

The first two statements of \ref{thm:AB-vl} follow easily from the fact that $i:A \to B$ induces a map of canonical Adams resolutions. The proof of the third statement will occupy most of Sections \ref{app:ABvl} and \ref{app:blocks}. In this proof we will rely on an alternative interpretation of $f_R$ from \cite{Akhil}. In order to recall this interpretation we begin by defining a natural filtration on the thick $\otimes$-ideal generated by $R$.

\begin{dfn}
  Given a set of spectra, $S$, the thick $\otimes$-ideal generated by $S$ consists of the smallest collection of spectra, $\text{Thick}^{\otimes}(S)$, closed under finite (co)limits and retracts, such that $X \otimes s \in \text{Thick}^{\otimes}(S)$ for all $s \in S$. We equip $\text{Thick}^{\otimes}(S)$ with the following filtration:
  \begin{itemize}
  \item $\text{Thick}^{\otimes}(S)_0 = \{0\}$,
  \item $\text{Thick}^{\otimes}(S)_1$ consists of retracts of spectra of the form $X \otimes s$ where $s \in S$,
  \item $\text{Thick}^{\otimes}(S)_n$ consists of retracts of extensions of objects of $\text{Thick}^{\otimes}(S)_{n-1}$ by objects of $\text{Thick}^{\otimes}(S)_1$.
  \end{itemize}
  We will only make use of this definition in the case where $S = \{R\}$.
\end{dfn}

\begin{rmk}
  For any $R$-module $M$ the unit map $M \to R \otimes M$ and the action map $R \otimes M \to M$ exhibit $M$ as a retract of $R \otimes M$, therefore $M \in \text{Thick}^\otimes(R)_1$.
\end{rmk}

The function $f_R$ can then be interpreted in terms of this filtration. 

\newcounter{fint}
\begin{prop}[{\cite[Definitions 2.28 and 3.26, and Proposition 3.28]{Akhil}}]
  \label{prop:f-interpret-I}\ 
  
  Let $I := \mathrm{fib}(\Ss \to R)$, then the following are equivalent:
  \begin{enumerate}
  \item $f_R(k) \leq n$,
  \item $\tau_{<k}\Ss^0 \in \textup{Thick}^{\otimes}(R)_n$,
  \item the map $I^{\otimes n} \to \Ss^0$ becomes null after tensoring with $\tau_{<k}\Ss^0$.
    \setcounter{fint}{\value{enumi}}    
  \end{enumerate}
\end{prop}

Sadly, none of these conditions are particularly convenient for the proof we have in mind. In order to remedy this we introduce two further equivalent conditions:

\begin{prop}
  \label{prop:f-interpret-II}
  The list of equivalent conditions from \Cref{prop:f-interpret-I} can be extended to include:
  \begin{enumerate}
    \setcounter{enumi}{\value{fint}}
  \item $\tau_{<k}\Ss^0$ is a retract of an object which has a length $n$ resolution by connective $R$-modules.
  \item $\tau_{<k}\Ss^0$ is a retract of an object which has a length $n$ resolution by connective induced $R$-modules.
  \end{enumerate}
\end{prop}

Before proving \Cref{prop:f-interpret-II} we set up some notation and conventions for manipulating finite resolutions of spectra.

\begin{dfn}
  A length $N$ resolution of a spectrum $X_0$ will consist of a diagram
  \begin{center}
    \begin{tikzcd}
        F_{N-1} \ar[r] & X_{N-2} \ar[r] \ar[d] & \cdots \ar[r] & X_1 \ar[r] \ar[d] & X_0 \ar[d] \\
        & F_{N-2} & & F_1 & F_0
    \end{tikzcd}
  \end{center}
  such that each $X_{j+1} \to X_j \to F_j$ is a cofiber sequence with the convention that $X_{N-1} = F_{N-1}$.
  In this situation we will say that we have a resolution of $X_0$ by $F_{N-1},\dots,F_0$.
\end{dfn}

\begin{ntn}
  We will adopt the following compact notation
  $$ [F_N, \dots, F_1, F_0; X] $$
  to express a resolution of $X$ by $F_{N},\dots,F_0$.
  It is important to note that this notation suppresses much of the data of a resolution.
\end{ntn}

\begin{wrn}
  Sometimes we will write $[\dots, F_1, F_0;X]$ for a resolution. Although this suggests an infinite-length resolution, in this appendix all resolutions will be finite length and this will simply be used to avoid specifying the length of a resolution. 
\end{wrn}

\begin{rmk}
  In the length two case the notation $[A, B; X]$ simply refers to a cofiber sequence
  $$ A \to X \to B. $$
\end{rmk}

\begin{proof}[Proof of \Cref{prop:f-interpret-II}]\
  
  $(5) \Rightarrow (4)$, clear.
  
  $(4) \Rightarrow (2)$,
  As remarked above every $R$-module is in $\text{Thick}^{\otimes}(R)_1$, therefore $\tau_{<k}\Ss^0$ is a retract of an $n$-fold extension of elements of $\text{Thick}^{\otimes}(R)_1$.
  
  $(3) \Rightarrow (5)$,
  Consider the following length $n+1$ resolution of $\Ss$,
  \begin{center}
    \begin{tikzcd}
      I^{\otimes n} \ar[r] &
       I^{\otimes (n-1)} \ar[r] \ar[d] &
      \cdots \ar[r] &
       I \ar[r] \ar[d] &
      \Ss \ar[d] \\
      &  R \otimes I^{\otimes (n-1)} & &
      R \otimes I &
      R
    \end{tikzcd}
  \end{center}
  From it we can produce a length $n$ resolution of $\text{cof}(I^{\otimes n} \to \Ss)$,
  \begin{center}
    \begin{tikzcd}
      R \otimes I^{\otimes (n-1)} \ar[r] &
      I^{\otimes (n-2)}/I^{\otimes n} \ar[r] \ar[d] &
      \cdots \ar[r] &
       I / I^{\otimes n} \ar[r] \ar[d] &
      \Ss / I^{\otimes n} \ar[d] \\
      &  R \otimes I^{\otimes (n-2)} & &
      R \otimes I &
      R
    \end{tikzcd}
  \end{center}
  Upon tensoring with $\tau_{<k}\Ss^0$ we obtain a length $n$ resolution of $(\tau_{<k}\Ss^0) \otimes (\Ss/I^{\otimes n})$ by connective, induced $R$-modules. Finally, by hypothesis,
  $$ (\tau_{<k}\Ss^0) \otimes (\Ss/I^{\otimes n}) \simeq (\tau_{<k}\Ss^0) \oplus (\tau_{<k}\Ss^0 \otimes \Sigma I^{\otimes n}). $$
\end{proof}

Using condition (4) we reduce the proof of the Theorem \ref{thm:AB-vl} to the following problem:
take a resolution of $X$ by $A$-modules and produce from it the shortest possible resolution of $X$ by $B$-modules. In order to provide a simple illustration of the methods we will use in the general case we first work the following example in detail.

\begin{qst*} Suppose that a spectrum $X$ sits in a cofiber sequence $C \to X \to D$ where $C,D$ are $\textup{BP}$-modules and $C,D,X \in \Sp_{[0,10]}$. What is the shortest resolution of $X$ by $\textup{BP}\langle 1 \rangle$-modules?
\end{qst*}

{\bf Strategy 1:}
We know that the map $\textup{BP} \to \textup{BP}\langle 1 \rangle$ is an equivalence after we apply $\tau_{<6}$, therefore any $\textup{BP}$ module in $\Sp_{[k,k+5]}$ is automatically a $\textup{BP}\langle 1 \rangle$ module.\footnote{A proof of this will appear in much greater generality in the next section} Knowing this trick we can break each of $C$ and $D$ into two $\textup{BP}\langle 1 \rangle$-modules and produce a new resolution of $X$ which uses 4 $\textup{BP}\langle 1 \rangle$-modules:
\begin{center}
  \begin{tikzcd}
    \tau_{[6,10]}C \ar[r] & C \ar[d] \ar[r] & F \ar[d] \ar[r] & X \ar[d] \\
    & \tau_{[0,5]}C & \tau_{[6,10]}D & \tau_{[0,5]}D
  \end{tikzcd}
\end{center}

This is a start, but it turns out we can do better.

{\bf Strategy 2:}
For our second approach we will start with a slightly modified version of the first resolution we produced:
\begin{center}
  \begin{tikzcd}
    \tau_{[5,10]}C \ar[r] & C \ar[d] \ar[r] & F \ar[d] \ar[r] & X \ar[d] \\
    & \tau_{[0,4]}C & \tau_{[6,10]}D & \tau_{[0,5]}D
  \end{tikzcd}
\end{center}
Now, we can expand this resolution into the diagram below where each square is cartesian.
\begin{center}
  \begin{tikzcd}
    \tau_{[5,10]}C \ar[r] \ar[d] &
    C \ar[r] \ar[d] &
    F \ar[r] \ar[d] &
    X \ar[d] \\
    0 \ar[r] &
    \tau_{[0,4]}C \ar[r] \ar[d] &
    G \ar[d] \ar[r] &
    H \ar[d] \\
    & 0 \ar[r] &
    \tau_{[6,10]}D \ar[r] \ar[d] &
    D \ar[d] \\
    & & 0 \ar[r] &
    \tau_{[0,5]}D 
  \end{tikzcd}
\end{center}

Notably the cofiber sequence which $G$ sits in is ``backwards''.
In fact, if we expand it a little bit
$$ \Sigma^{-1}\tau_{[6,10]}D \to \tau_{[0,4]}C \to G \to \tau_{[6,10]}D $$
we see that the attaching map must be zero for connectivity reasons and therefore $G \simeq \tau_{[0,4]}C \oplus \tau_{[6,10]}D$.
As a direct sum of $\textup{BP}\langle 1 \rangle$-modules this is in fact a $\textup{BP}\langle 1 \rangle $-module as well. Thus, we can produce the following length $3$ resolution of $X$
\begin{center}
  \begin{tikzcd}
    \tau_{[5,10]}C \ar[r] & F \ar[d] \ar[r] & X \ar[d] \\
    & \tau_{[0,4]}C \oplus \tau_{[6,10]}D & \tau_{[0,5]}D
  \end{tikzcd}
\end{center}

%% \begin{rmk}
%%   The first step of strategy 2 was to slightly modify the resolution from strategy 1. Although $G$ does not split without this modification it is still a $\textup{BP}\langle 1 \rangle$-module. This more complicated argument would give a sharper bound on $f_B$ in \Cref{thm:AB-vl}.
%% \end{rmk}

Both strategies had four main steps:
\begin{enumerate}
\item start with a resolution by bounded $A$-modules,
\item construct a new resolution from the given one,
\item count the length of the new resolution.
\item show that the new resolution is in fact a resolution by $B$-modules,
\end{enumerate}
In the proof of Theorem \ref{thm:AB-vl} each of these elements will be replaced by a lemma (whose proofs we defer to the next section).
% The key step is the construction of a new resolution and ultimately this is no more than a systematic elaboration of strategy 2 above.

\begin{lem}
  \label{lemm:resolution-prep}
  Given a resolution $[F_N,\dots,F_0; X]$ such that each $F_j$ is an $A$-module and all the $F_j$ are connective, there exists another resolution $[F_N',\dots,F_0'; \tau_{\leq M-1}X]$ such that each $F_j'$ is an $A$-module in $\Sp_{[0,M]}$.
\end{lem}

\begin{lem}
  \label{lemm:new-resolution}
  Given a resolution $[F_N, \dots, F_0; X]$ where each $F_j$ lives in $\Sp_{[0,K]}$
  we can construct another resolution
  $$ \left[ \dots, \left( \bigoplus_{0 \leq i \leq j} \tau_{[(j-i)m-i,(j-i+1)m-i)}F_i \right) , \dots, \left( \tau_{[-1,m-1)}F_1  \oplus  \tau_{[m,2m)}F_0 \right) , \left( \tau_{[0,m)}F_0 \right) ; X \right] $$
\end{lem}

\begin{rmk}
  Since \Cref{lemm:new-resolution} is more complicated than the others we pause here to note that in the case where $N=0$ this lemma just produces the $m$-speed Postnikov tower for $X$.
  In general, the lemma takes Postnikov-type towers for each $F_j$ and shuffles them together.
  Ultimately this is no more than a careful elaboration on the manipulations used in strategy 2 above.
\end{rmk}

\begin{lem}
  \label{lemm:length-count}
  The resolution produced by the Lemma \ref{lemm:new-resolution} has length,
  $$ 1 + N + \left\lfloor \frac{K + N}{m} \right\rfloor. $$
\end{lem}

\begin{lem}
  \label{lemm:res-by-B-mod}
  Given a ring map $A \to B$ which becomes an equivalence after applying $\tau_{<m}$ any $A$-module in $\Sp_{[a,a+m)}$ can be given the structure of a $B$-module.
\end{lem}

\begin{proof}[Proof of Theorem \ref{thm:AB-vl}]
  %% invoke the A resolution
  Let $N+1 = f_A(k)$, then by Proposition \ref{prop:f-interpret-II}(4) there exists a $Y$ such that
  \begin{enumerate}
  \item $\tau_{<k}\Ss^0$ is a retract of $Y$ and
  \item $Y$ has a resolution $[F_{N},\dots,F_0; Y]$ by connective $A$-modules.
  \end{enumerate}
  Next we apply Lemma \ref{lemm:resolution-prep} to obtain a resolution $[G_{N},\dots,G_0; \tau_{<k}Y]$ where each $G_j$ is an $A$-module in $\Sp_{[0,k]}$.

  %% obtain the new resolution
  At this point we apply Lemma \ref{lemm:new-resolution} to
  $ [G_{N},\dots,G_0; \tau_{<k}Y] $ with $K = k$, $m = m$ to obtain a new resolution
  $ [\dots,H_1,H_0; \tau_{<k}Y] $.
  %% prove it's a B resolution
  Each of the $H_j$ is a direct sum of finitely many terms of the form $\tau_{[a,a+m)}G_i$.
  By Lemma \ref{lemm:res-by-B-mod} each of these terms is then a $B$-module,
  thus $H_j$ is a $B$-module as well.
  Finally, we note that $\tau_{<k}\Ss^0$ is a retract of $\tau_{<k}Y$, therefore by Proposition \ref{prop:f-interpret-II}(4) $f_B(k)$ is bounded by the length of the resolution we have produced and Lemma \ref{lemm:length-count} lets us conclude that
  $$ f_B(k) \leq 1 + N + \left\lfloor \frac{k + N}{m} \right\rfloor. $$
\end{proof}

%%% Local Variables:
%%% mode: latex
%%% TeX-master: "Main"
%%% eval: (visual-line-mode t)
%%% eval: (auto-fill-mode 0)
%%% End:
